%=====================================================================
%  Electrical Design of the Rail-Guided TEC PV-Cleaning Robot
%  and its Atmospheric-Water-Generator (AWG) Dock
%
%  Project : Graduation Project (REE598)
%  Site    : New Administrative Capital, Egypt (residential dock)
%  Output  : .tex  |  Units: SI (siunitx)  |  Citations: IEEE (IEEEtran)
%
%  Numerical inputs are taken from the verified design engine outputs:
%    output/rail_guided_design.json
%    output/tec_robot_results.json
%    output/nac_tec_robot_results.json
%  Part numbers and prices: MARKET_PARTS_BUYLIST.md
%=====================================================================
\documentclass[11pt,a4paper]{article}

% --- Core --------------------------------------------------------------
\usepackage[utf8]{inputenc}
\usepackage[T1]{fontenc}
\usepackage{lmodern}
\usepackage[english]{babel}

% --- Mathematics & units ----------------------------------------------
\usepackage{amsmath,amssymb}
\usepackage{mathtools}
\usepackage{textcomp}
\usepackage{siunitx}
\sisetup{per-mode=symbol, detect-weight=true, detect-family=true}
\DeclareSIUnit\year{yr}
\DeclareSIUnit\rpm{rpm}

% --- Figures & tables (no colortbl: MiKTeX colortbl/array bug) ---------
\usepackage{graphicx}
\usepackage{float}
\usepackage{booktabs}
\usepackage{array}
\usepackage{caption}
\usepackage{longtable}

% --- Layout ------------------------------------------------------------
\usepackage{geometry}
\geometry{margin=2.3cm}
\usepackage{setspace}
\onehalfspacing
\usepackage{cite}
\usepackage{hyperref}
\hypersetup{colorlinks=true,linkcolor=blue,citecolor=blue,urlcolor=blue}

\newcommand{\degC}{\si{\degreeCelsius}}

\title{\bfseries Electrical Design of the Rail-Guided\\
Thermoelectric PV-Cleaning Robot and its\\
Atmospheric-Water-Generator Dock}
\author{Rahma El-Hamouly \\ Graduation Project --- REE598}
\date{\today}

\begin{document}
\maketitle

\begin{abstract}
\noindent
This document derives the complete electrical design of a two-part PV-cleaning
system: a \SI{8}{\kilo\gram} rail-guided carriage that brushes and sprays a
photovoltaic array, and a stationary atmospheric-water-generator (AWG) dock that
condenses the carriage's cleaning water from air using thermoelectric coolers
(TECs). The carriage runs on a \SI{12}{\volt} lithium-iron-phosphate (LFP)
battery; the dock runs a \num{38}-tile TEC bank from PV with maximum power-point
tracking (MPPT) and storage. The design follows a seven-step workflow: power
budget, supply, motor sizing and drive, energy storage, wire and fuse sizing,
protection, and TEC electrical design. Two \num{1}:\num{71} GA25-370 gearmotors
drive the rack and pinion through one right-sized Cytron MDD3A dual driver, and
the carriage battery is held to \SI{10}{\ampere\hour} so the carriage closes
under its \SI{8}{\kilo\gram} mass limit while still giving about five full-array
cleaning sessions per charge. The headline finding concerns the dock: running
\num{38} TEC tiles at
\SI{6}{\ampere} each on a \SI{12}{\volt} bus draws roughly \SI{228}{\ampere},
which forces \SI{36}{\milli\meter\squared} copper and dumps about \SI{14}{\percent}
of the bank power as $I^2R$ loss. Raising the TEC bus to \SI{48}{\volt} cuts the
current to \SI{63}{\ampere}, the conductor to \SI{2.3}{\milli\meter\squared}, and
the cable loss to under \SI{1}{\percent}, and roughly halves the MPPT cost. The
\SI{12}{\volt} carriage bus is kept as-is because its power is low.
\end{abstract}

\tableofcontents
\bigskip

%=====================================================================
\section{System architecture and scope}
\label{sec:arch}

The system splits into two electrically independent subsystems that share no
power bus.

\textbf{(A) The carriage} is the rail-guided robot. It carries a rack-and-pinion
drive (two driven rails), a full-width brush roller, a low-pressure water pump, a
small sensor suite, and a two-tier compute stack (an ESP32 real-time controller
and a Raspberry~Pi~4 host). It is powered by a \SI{12}{\volt} LFP battery and
draws at most a few tens of watts. The carriage is the \emph{actuator} side of
the design and is dominated by motion and spray, not by computation.

\textbf{(B) The dock} is the stationary AWG. It runs \num{38} TEC1-12706 tiles
against finned hot-side heat sinks with axial fans, plus a microcontroller for
the night-harvest strategy. It is powered by a PV array through an MPPT charge
controller and a battery bank, or, in the recommended residential case,
grid-coupled to the house's existing \SI{9.9}{\kilo\watt_{p}} roof array. The
dock is the \emph{thermal} side and is dominated by TEC electrical power, which
is where the bus-voltage decision lives.

All numerical inputs come from the verified design engine outputs
(\texttt{rail\_guided\_design.json}, \texttt{tec\_robot\_results.json},
\texttt{nac\_tec\_robot\_results.json}); component choices and prices come from
the market buy-list. Datasheet citations are given only for parts that have a
real datasheet entry.

%=====================================================================
\section{Step 1 --- Power budget}
\label{sec:budget}

The power budget is built first, before any component is selected, so that every
later choice has a quantified demand behind it.

\subsection{Carriage power budget}

The carriage operates in a duty-cycled pattern: it sweeps one panel width in
\SI{66.7}{\second} (a \SI{33.3}{\second} brushing-and-spraying down-pass and a
\SI{33.3}{\second} return), then repositions to the next module. Table~\ref{tab:carriage-budget}
lists the electrical loads. The motor currents are derived in
Section~\ref{sec:motors} and are quoted here for completeness.

\begin{table}[H]
\centering
\caption{Carriage electrical power budget. Duty cycle is the fraction of a
cleaning session each load is energised.}
\label{tab:carriage-budget}
\small
\begin{tabular}{l S[table-format=3.1] S[table-format=3.0] S[table-format=2.1]}
\toprule
Load & {Peak power (\si{\watt})} & {Duty (\si{\percent})} & {Avg.\ (\si{\watt})} \\
\midrule
Drive motors ($2\times$ GA25-370 1:71) & 19.3 & 100 & 19.3 \\
Brush motor (GA25-370, steady)    & 1.6  & 50  & 0.8 \\
Water pump (diaphragm, spraying)  & 45.0 & 50  & 22.5 \\
ESP32 controller                  & 1.0  & 100 & 1.0 \\
Raspberry Pi 4 (active)           & 7.0  & 100 & 7.0 \\
Sensors (DHT22, $2\times$DS18B20, INA219, MPU-6050, limits) & 0.4 & 100 & 0.4 \\
\midrule
\textbf{Total (motion + spray)}   & \textbf{74.3} & {} & \textbf{51.0} \\
\bottomrule
\end{tabular}
\end{table}

The water pump dominates both the instantaneous draw and, with the lighter
\num{1}:\num{71} drive motors, the time-averaged motion power. The Raspberry~Pi
can be gated off
entirely when the robot is parked and only the dock is harvesting, removing
\SI{5}{\watt} of standby; that gating is assumed in the energy-storage sizing of
Section~\ref{sec:storage}.

The energy consumed per cleaning session, summing each load over its active time
$t_i$, is
\begin{equation}
  E_{\text{cycle}} = \sum_i P_i\, t_i .
  \label{eq:ecycle}
\end{equation}
Here $P_i$ is the power of load $i$ and $t_i$ its energised time in the session.
For one panel sweep the down-pass (drive, brush, pump, logic) consumes
\SI{0.69}{\watt\hour} and the return (drive and logic only)
\SI{0.26}{\watt\hour}, giving \SI{0.94}{\watt\hour} per panel. A residential
session covers the \num{18}-module array, so

\begin{equation}
  E_{\text{session}} = 18 \times \SI{0.94}{\watt\hour} + E_{\text{idle}}
  \approx \SI{17.3}{\watt\hour},
  \label{eq:esession}
\end{equation}
where $E_{\text{idle}}\approx\SI{0.3}{\watt\hour}$ accounts for repositioning
dwell at \SI{6}{\watt} idle power. This figure sizes the carriage battery. It is
lower than the earlier estimate because the \num{1}:\num{71} drive motors draw
\SI{19.3}{\watt} rather than the \SI{36.5}{\watt} of the faster variant.

\subsection{Dock power budget}

The dock budget is set by the TEC bank, not by control electronics. From the NAC
design at full self-supply (\num{38} tiles), the daily energy split is given in
Table~\ref{tab:dock-budget}.

\begin{table}[H]
\centering
\caption{Dock daily energy budget at \num{38} tiles, full atmospheric
self-supply (from \texttt{nac\_tec\_robot\_results.json}, $f=1.0$).}
\label{tab:dock-budget}
\small
\begin{tabular}{l S[table-format=4.0] S[table-format=2.1]}
\toprule
Load & {Daily energy (\si{\watt\hour})} & {Share (\si{\percent})} \\
\midrule
TEC tiles ($38\times$)         & 4206 & 65.0 \\
Hot-side fans ($20\times\,\SI{2.5}{\watt}$) & 1080 & 16.7 \\
Dock MCU + telemetry           & 36   & 0.6 \\
Cleaning-water pumping (shared) & 1   & 0.0 \\
\midrule
\textbf{Total}                 & \textbf{6470} & {} \\
\bottomrule
\end{tabular}
\end{table}

The TEC tiles draw \SI{65}{\percent} of the dock energy. Everything downstream of
this table follows from how those tiles are wired and at what bus voltage, which
is the subject of Section~\ref{sec:tec}.

%=====================================================================
\section{Step 2 --- Power supply}
\label{sec:supply}

\subsection{Carriage supply}

The carriage is battery-powered at \SI{12}{\volt} nominal. The logic rails
(\SI{5}{\volt} for the Pi and \SI{3.3}{\volt} for the ESP32 and sensors) come
from a switching buck regulator. For a switcher the loss is
\begin{equation}
  P_{\text{loss}} = (1-\eta_{sw})\,P_{\text{out}},
  \label{eq:buckloss}
\end{equation}
with $\eta_{sw}\approx\num{0.9}$ for an LM2596-class module. The \SI{5}{\volt}
rail carries the Pi (\SI{7}{\watt}) plus sensors, so $P_{\text{out}}\approx
\SI{8}{\watt}$ and $P_{\text{loss}}\approx\SI{0.9}{\watt}$, dissipated in the
regulator. A \SI{3}{\ampere} buck module covers this with margin.

\subsection{Dock supply from PV}

The dock's daily PV energy is
\begin{equation}
  E_{\text{day}} = P_{\text{STC}}\cdot PR \cdot H_{\text{peak}},
  \label{eq:eday}
\end{equation}
where $P_{\text{STC}}$ is the array rating, $PR$ the performance ratio
($\approx\num{0.80}$), and $H_{\text{peak}}$ the peak sun hours
(\SI{5.71}{\hour} at the NAC site). To deliver the \SI{6470}{\watt\hour\per\day}
of Table~\ref{tab:dock-budget} after storage round-trip losses, the design engine
sizes the array at $P_{\text{STC}}\approx\SI{1416}{W_{p}}$ (three
\SI{550}{W_{p}} modules). The grid-coupled alternative skips the dedicated PV,
MPPT, and bank, drawing the dock's night load from the house's existing
\SI{9.9}{\kilo\watt_{p}} array; that option leaves only about
\SI{14000}{EGP} of marginal AWG hardware and is the recommended residential
deployment.

%=====================================================================
\section{Step 3 --- Motor sizing and drive}
\label{sec:motors}

\subsection{Drive motors}

The rail drive must overcome the total resistive force on the \SI{24}{\degree}
rail. From \texttt{rail\_guided\_design.json}, gravity along the slope
contributes \SI{31.9}{\newton}, rolling resistance \SI{3.6}{\newton}, and
brushing drag \SI{1.7}{\newton}, for a total
\begin{equation}
  F_{\text{drive}} = F_{g} + F_{\text{roll}} + F_{\text{clean}}
  = \SI{37.2}{\newton}.
  \label{eq:fdrive}
\end{equation}
At a \SI{30}{\milli\meter} pinion (radius $r=\SI{15}{\milli\meter}$) the traction
torque is
\begin{equation}
  T_{\text{trac}} = F_{\text{drive}}\, r = \SI{37.2}{\newton}\times\SI{0.015}{\meter}
  = \SI{0.838}{\newton\meter}.
  \label{eq:ttrac}
\end{equation}
The design uses two driven rails, so the per-motor load is
$T_{\text{load}} = \SI{0.838}{\newton\meter}/2 = \SI{0.42}{\newton\meter}$.

The selected motor is the GA25-370 \SI{12}{\volt} geared motor with a quadrature
encoder, in the \num{1}:\num{71} (``\SI{60}{\rpm}'') reduction variant. Its
datasheet figures at \SI{12}{\volt} are \SI{0.78}{\newton\meter} stall torque,
\SI{60}{\rpm} no-load, \SI{1.5}{\ampere} stall current, and \SI{0.42}{\ampere}
rated current. The higher reduction is a deliberate choice over the
\SI{250}{\rpm} variant: it keeps the output near the cleaning speed (so the PWM
duty sits in a controllable range, not pinned near \num{1}) and it raises the
reflected gearbox friction that resists back-driving on the slope, which is the
slide problem Section~\ref{sec:hold} addresses.

The torque constant referred to the gearbox output is
\begin{equation}
  k_t = \frac{T_{\text{stall}}}{I_{\text{stall}}}
  = \frac{\SI{0.78}{\newton\meter}}{\SI{1.5}{\ampere}}
  = \SI{0.520}{\newton\meter\per\ampere},
  \label{eq:kt}
\end{equation}
where $T_{\text{stall}}$ is the stall torque, $I_{\text{stall}}=\SI{1.5}{\ampere}$
the stall current, and $k_t$ the torque constant. The motor follows the standard
DC model
\begin{equation}
  V = k_e\,\omega + I R_a, \qquad T = k_t I,
  \label{eq:dcmodel}
\end{equation}
with $V$ the terminal voltage, $k_e$ the back-EMF constant, $\omega$ the angular
speed, $I$ the armature current, $R_a$ the armature resistance, and $T$ the
output torque. Applying $T = k_t I$ at the operating load,
\begin{equation}
  I_{\text{op}} = \frac{T_{\text{load}}}{k_t}
  = \frac{\SI{0.42}{\newton\meter}}{\SI{0.520}{\newton\meter\per\ampere}}
  = \SI{0.81}{\ampere}\ \text{per motor.}
  \label{eq:iop_drive}
\end{equation}
Two motors draw \SI{1.61}{\ampere} together at \SI{12}{\volt}, or
\SI{19.3}{\watt} electrical. The mechanical output is
$P_{\text{shaft}} = T_{\text{trac}}\,\omega = \SI{2.2}{\watt}$ at
\SI{38}{\rpm}, so most of the electrical input is gearbox and motor loss at this
low load, which is expected for a small geared motor run well below its rated
torque.

The torque margin is $\SI{0.78}{\newton\meter}/\SI{0.42}{\newton\meter}=
\num{1.86}\times$ per motor, and the speed margin is
$\SI{60}{\rpm}/\SI{38}{\rpm}\approx\num{1.57}\times$. Both clear the
\num{1.5} selection floor.

\subsection{Holding the carriage on the tilt}
\label{sec:hold}

The earlier free-wheel crawler slid down the glass when unpowered; the rail drive
must not repeat that. Gravity along the \SI{24}{\degree} slope is
\SI{31.9}{\newton}, which referred back through the \SI{15}{\milli\meter} pinion
is a back-drive torque of
\begin{equation}
  T_{\text{back}} = F_{g}\,r = \SI{31.9}{\newton}\times\SI{0.015}{\meter}
  = \SI{0.48}{\newton\meter},
  \label{eq:tback}
\end{equation}
or \SI{0.24}{\newton\meter} per rail. A GA25-370 spur or planetary gearbox is not
self-locking, so the \num{1}:\num{71} reduction alone cannot be trusted to hold
this: the motor would creep. The committed hold strategy is the motor driver's
brake mode, which shorts the two motor terminals at park so the back-EMF of any
creep develops a retarding current, plus a small sprung detent that drops into the
rack at the docked position for a positive mechanical stop with no standing
current. Where a brake-free hold is wanted, the JGY-370 worm gearmotor is the
drop-in alternative: a worm drive is geometrically self-locking, so it holds the
carriage with the supply off at the cost of some efficiency and a narrower ratio
choice. The buy-list carries both options.

\subsection{Brush motor}

The full-width \SI{0.96}{\meter} roller has \SI{1.73}{\newton} of drag and a
required torque of \SI{0.065}{\newton\meter} at \SI{100}{\rpm}, for about
\SI{0.5}{\watt} mechanical. A GA25-370 \SI{100}{\rpm} variant
($k_t\approx\SI{0.49}{\newton\meter\per\ampere}$) gives a steady current of
$\SI{0.065}{\newton\meter}/k_t=\SI{0.13}{\ampere}$, a large margin against steady
drag. The real sizing driver is not steady drag but the start-up transient: the
roller and its water film must be spun up from rest, and the inrush approaches the
stall current of \SI{1.6}{\ampere} for a fraction of a second. The switch and fuse
are sized to that inrush, not to the \SI{0.13}{\ampere} steady value.

\subsection{Water pump}

The pump is a \SI{12}{\volt} diaphragm unit in the \SIrange{3.6}{4.0}{\liter\per\minute}
class for low-pressure spray. Measured current for this class while spraying is
\SIrange{3}{4}{\ampere}, so \SI{45}{\watt} is taken as the design point. The pump
runs only during the down-pass, a \SI{50}{\percent} duty over the sweep, which is
why its time-averaged contribution in Table~\ref{tab:carriage-budget} is half its
peak.

\subsection{Driver and switch selection}
\label{sec:driver}

The two drive motors need reversal (up and down the rail), so they take a
proper H-bridge; the brush and pump run in one direction only, so each takes a
single low-side switch. The driver selection rules are: continuous rating
$\geq 1.5\,I_{\text{op}}$, peak rating $\geq I_{\text{stall}}$, and voltage
rating $\geq 1.3\,V_{\text{supply}}$.

A \SI{43}{\ampere} BTS7960 module against a motor whose stall current is
\SI{1.5}{\ampere} is roughly fifty times over-rated, which wastes board space,
cost, and quiescent current. The committed driver is instead the Cytron MDD3A
\cite{cytronmdd3a}, a dual-channel \SI{3}{\ampere}-continuous,
\SI{5}{\ampere}-peak, \SIrange{5}{30}{\volt} H-bridge, so one module drives both
rail motors. Checking it against the drive motor,
\begin{align}
  \SI{3}{\ampere} &\geq 1.5\times\SI{0.81}{\ampere} = \SI{1.21}{\ampere}, \\
  \SI{5}{\ampere} &\geq I_{\text{stall}} = \SI{1.5}{\ampere}, \\
  \SI{30}{\volt}  &\geq 1.3\times\SI{12}{\volt} = \SI{15.6}{\volt}.
  \label{eq:mdd3a}
\end{align}
The continuous headroom is \num{3.7}$\times$ the operating current, which is the
sizing intent: enough to ride the inrush, not so much that the part is dead
weight. The MDD3A also supports a brake mode, which Section~\ref{sec:hold} uses
to hold the carriage on the slope. Note that the common small dual drivers fail
the voltage rule here: the TB6612FNG ceiling is \SI{13.5}{\volt} and the DRV8833
is \SI{10.8}{\volt}, both below the \SI{15.6}{\volt} the \num{1.3}$\times$ rule
asks for on a \SI{12}{\volt} bus, so the MDD3A is the smallest part that still
clears all three checks.

The brush and pump each switch with an IRLZ44N logic-level N-MOSFET
\cite{irlz44nds} on the low side, gated directly from an ESP32 pin through a gate
resistor, with a flyback diode across each inductive load. The IRLZ44N
(\SI{55}{\volt}, \SI{47}{\ampere}) is far above the brush \SI{1.6}{\ampere}
inrush and the pump \SI{3.75}{\ampere}, but it is the cheapest robust
logic-level part on the shelf and avoids a second H-bridge for loads that never
reverse.

\subsection{PWM speed control}

Speed is set by PWM duty
\begin{equation}
  D = \frac{V_{\text{desired}}}{V_{\text{supply}}},
  \label{eq:duty}
\end{equation}
where $V_{\text{desired}}$ is the average voltage needed for the target speed and
$V_{\text{supply}}=\SI{12}{\volt}$. The drive target of \SI{38}{\rpm} from the
\SI{60}{\rpm} no-load motor needs $D\approx\num{0.64}$, comfortably inside the
controllable range with droop headroom; the \SI{250}{\rpm} variant would have
pinned the duty near \num{0.15}, where load droop is harder to regulate. The
brush runs near its rated speed at $D\approx\num{1.0}$. The PWM frequency is set
to
$f_{\text{PWM}}\geq\SI{20}{\kilo\hertz}$ to stay above the audible band. The
resulting current ripple is
\begin{equation}
  \Delta I = \frac{V_{\text{supply}}\,D\,(1-D)}{L\,f_{\text{PWM}}},
  \label{eq:ripple}
\end{equation}
with $L$ the motor inductance ($\approx\SI{1}{\milli\henry}$ for this motor
class). At the worst-case $D=\num{0.5}$,
$\Delta I = \SI{12}{\volt}\times\num{0.25}/(\SI{1e-3}{\henry}\times\SI{20e3}{\hertz})
= \SI{0.15}{\ampere}$, about \SI{10}{\percent} of the operating current and well
within the driver's capability.

%=====================================================================
\section{Step 4 --- Energy storage (carriage battery)}
\label{sec:storage}

The carriage battery capacity from a single-cycle energy budget is
\begin{equation}
  C_{\text{battery}} = \frac{E_{\text{total}}}{V_{\text{nom}}\cdot DOD\cdot\eta_{c}},
  \label{eq:batt}
\end{equation}
where $E_{\text{total}}$ is the energy per cleaning session, $V_{\text{nom}}=
\SI{12}{\volt}$, $DOD=\num{0.8}$ for LFP, and $\eta_{c}=\num{0.95}$ the charge
efficiency. Substituting $E_{\text{total}}=\SI{17.3}{\watt\hour}$ from
Eq.~\eqref{eq:esession},
\begin{equation}
  C_{\text{battery}} = \frac{\SI{17.3}{\watt\hour}}{\SI{12}{\volt}\times\num{0.8}\times\num{0.95}}
  = \SI{1.90}{\ampere\hour}.
  \label{eq:batt_num}
\end{equation}
With a \SI{20}{\percent} design margin, $C_{\text{actual}}=\num{1.2}\times
\SI{1.90}{\ampere\hour}=\SI{2.3}{\ampere\hour}$. A single cleaning session is
therefore trivially small for any practical battery.

The pack is sized by two competing constraints, not by one session. Autonomy
wants a large pack so the robot cleans for several days between charges; the
carriage mass budget wants a small one. Section~\ref{sec:weight} shows that an
honest \SI{20}{\ampere\hour} LFP pack weighs about \SI{2.3}{\kilo\gram} and, with
the dock condenser already accounted for, would carry the carriage past its
\SI{8}{\kilo\gram} limit. The design therefore commits to a \SI{12}{\volt}
\SI{10}{\ampere\hour} LFP pack (about \SI{1.3}{\kilo\gram}) with a
\SI{15}{\ampere} BMS, and relies on a dock top-up between cleaning runs rather
than on a single large pack. Its usable energy is
\begin{equation}
  E_{\text{usable}} = C_{\text{pack}}\,V_{\text{nom}}\,DOD\,\eta_{c}
  = \SI{10}{\ampere\hour}\times\SI{12}{\volt}\times\num{0.8}\times\num{0.95}
  = \SI{91}{\watt\hour}.
  \label{eq:eusable}
\end{equation}
Dividing by the session energy gives the autonomy:
\begin{equation}
  n_{\text{sessions}} = \frac{E_{\text{usable}}}{E_{\text{session}}}
  = \frac{\SI{91}{\watt\hour}}{\SI{17.3}{\watt\hour}} \approx 5
  \ \text{full-array cleanings per charge.}
  \label{eq:autonomy}
\end{equation}
At the residential cleaning interval of \num{3.5} days, five sessions is roughly
\num{18} days of autonomy, after which the carriage tops up at the dock. The peak
carriage current while spraying is
\begin{equation}
  I_{\text{peak}} = \frac{P_{\text{peak}}}{V_{\text{nom}}}
  = \frac{\SI{74.3}{\watt}}{\SI{12}{\volt}} = \SI{6.2}{\ampere},
  \label{eq:ipeak}
\end{equation}
which is \SI{0.62}{C} on the \SI{10}{\ampere\hour} pack, below the
\SI{15}{\ampere} (\SI{1.5}{C}) BMS limit. The continuous runtime at the session
average power of \SI{51}{\watt} is about \SI{1.8}{\hour}, longer than the
\SI{20}{\minute} a full-array session takes. Coulomb counting through the INA219
tracks state of charge as
\begin{equation}
  SoC(t) = SoC_0 - \frac{1}{C_{\text{rated}}}\int_0^{t} I(\tau)\,d\tau,
  \label{eq:soc}
\end{equation}
with $SoC_0$ the initial charge, $C_{\text{rated}}$ the rated capacity, and
$I(\tau)$ the instantaneous current.

%=====================================================================
\section{Step 7 (placed early) --- TEC electrical design and the dock bus voltage}
\label{sec:tec}

The TEC electrical model is taken out of workflow order and placed before wire
sizing because the dock bus-voltage decision depends on it, and that decision is
the headline of this design.

\subsection{Per-tile operating point}

Each TEC1-12706 tile is modelled by the standard thermoelectric relations
\cite{tec12706ds,riffat2003peltier}:
\begin{align}
  Q_C &= \alpha T_C I - \tfrac{1}{2} I^2 R - K\,\Delta T, \label{eq:qc}\\
  Q_H &= \alpha T_H I + \tfrac{1}{2} I^2 R - K\,\Delta T, \label{eq:qh}\\
  V   &= \alpha\,\Delta T + I R, \qquad P_{\text{tile}} = V I, \label{eq:vtile}
\end{align}
where $Q_C$ is the heat absorbed at the cold side (the condensation duty), $Q_H$
the heat rejected at the hot side, $\alpha=\SI{0.0513}{\volt\per\kelvin}$ the
Seebeck coefficient, $R=\SI{2.002}{\ohm}$ the module resistance,
$K=\SI{0.546}{\watt\per\kelvin}$ the thermal conductance, $T_C$ the cold-side
temperature, $T_H$ the hot-side temperature, $\Delta T = T_H - T_C$, and $I$ the
tile current. These parameter values are from
\texttt{tec\_robot\_results.json}.

The operating point is set for condensation: a cold plate at \SI{10}{\degreeCelsius}
($T_C=\SI{283}{\kelvin}$) against a hot side near ambient
($\Delta T=\SI{25}{\kelvin}$), at the datasheet maximum current
$I=\SI{6}{\ampere}$ (within $I_{\max}=\SI{6}{\ampere}$,
$V_{\max}=\SI{15.4}{\volt}$). Substituting into Eqs.~\eqref{eq:qc}
and~\eqref{eq:vtile}:
\begin{align}
  V_{\text{tile}} &= \SI{0.0513}{\volt\per\kelvin}\times\SI{25}{\kelvin}
    + \SI{6}{\ampere}\times\SI{2.002}{\ohm} = \SI{13.3}{\volt}, \\
  P_{\text{tile}} &= \SI{13.3}{\volt}\times\SI{6}{\ampere} = \SI{79.8}{\watt}, \\
  Q_C &= \SI{37.4}{\watt}.
  \label{eq:tile_op}
\end{align}
The terminal voltage of \SI{13.3}{\volt} is below $V_{\max}$ and the tile power
of \SI{79.8}{\watt} is within the module rating, so the operating point is valid.

\subsection{COP: instantaneous versus system}

The instantaneous coefficient of performance at this operating point is
\begin{equation}
  COP_{\text{TEC}} = \frac{Q_C}{P_{\text{tile}}}
  = \frac{\SI{37.4}{\watt}}{\SI{79.8}{\watt}} = \num{0.47}.
  \label{eq:cop}
\end{equation}
This is the honest instantaneous figure, and it is low, as expected for a
single-stage TEC pumping against \SI{25}{\kelvin}. The system COP quoted elsewhere
in the project (\num{1.9}) is a different and larger number: it is the
\emph{seasonal} effective COP of the smart night-harvest strategy, which only
runs the tiles when night air is cool and humid, so the average $\Delta T$ over
the harvest hours is much smaller than the worst-case design point. The per-tile
mean COP across the harvest year is \num{1.44}; the NAC night-strategy effective
COP is \num{1.9}. Both the instantaneous \num{0.47} and the system \num{1.9} are
correct for their own definitions, and the design must not confuse them: the
\num{0.47} sets the electrical bus current, while the \num{1.9} sets the daily
energy and PV sizing.

\subsection{The 12\,V bus problem}

At the design operating point, \num{38} tiles wired in parallel on a single
\SI{12}{\volt} bus would each draw \SI{6}{\ampere}, for a bus current of
\begin{equation}
  I_{\text{bus,12V}} = 38\times\SI{6}{\ampere} = \SI{228}{\ampere}.
  \label{eq:ibus12}
\end{equation}
The total electrical power is
$P_{\text{array}} = 38\times\SI{79.8}{\watt} = \SI{3032}{\watt}$. A current of
\SI{228}{\ampere} on a \SI{12}{\volt} bus is impractical: it forces very heavy
copper, a huge MPPT current rating, and large $I^2R$ losses. The buy-list already
flagged that the \SI{12}{\volt} bus oversizes the MPPT controllers; the next
subsection quantifies the penalty and the fix.

\subsection{Bus-voltage trade: 12\,V vs 24\,V vs 48\,V}

The same array power $P_{\text{array}}=\SI{3032}{\watt}$ can be delivered at a
higher bus voltage by wiring the tiles in series strings: pairs for a
\SI{24}{\volt} bus (two tiles at \SI{13.3}{\volt} make \SI{26.6}{\volt}), or
quads for a \SI{48}{\volt} bus (four tiles make \SI{53.2}{\volt}). The
\num{38}-tile bank does not divide evenly into quads, so the wiring is nine
four-tile strings (\num{36} tiles) plus one two-tile string, giving ten strings
in parallel; the odd pair runs at the same per-tile \SI{6}{\ampere} string
current and its own switch and fuse. The string current stays at the per-tile
\SI{6}{\ampere}; only the number of parallel strings, and hence the bus current,
falls. The bus current scales inversely with voltage:
\begin{equation}
  I_{\text{bus}} = \frac{P_{\text{array}}}{V_{\text{bus}}}.
  \label{eq:ibus}
\end{equation}

The minimum conductor cross-section for a \SI{3}{\percent} voltage-drop budget is
\begin{equation}
  A_{\min} = \frac{I_{\text{bus}}\,L\,\rho_{\text{Cu}}}{V_{\text{drop,max}}},
  \qquad V_{\text{drop,max}} = \num{0.03}\,V_{\text{bus}},
  \label{eq:amin}
\end{equation}
with $\rho_{\text{Cu}}=\SI{1.72e-8}{\ohm\meter}$, $L=\SI{3}{\meter}$ the one-way
run, and $V_{\text{drop,max}}$ the allowed drop. The cable $I^2R$ loss for a
fixed \SI{16}{\milli\meter\squared} conductor over a round-trip $2L$ is
\begin{equation}
  P_{\text{loss}} = I_{\text{bus}}^2 \, R_{\text{wire}},
  \qquad R_{\text{wire}} = \frac{\rho_{\text{Cu}}\,(2L)}{A}.
  \label{eq:i2r}
\end{equation}
Table~\ref{tab:busvolt} collects the result. Both the conductor and the loss
scale with $1/V_{\text{bus}}^2$ at fixed power, because halving the current both
halves the required area for a given drop and quarters the $I^2R$ loss.

\begin{table}[H]
\centering
\caption{Dock TEC bus at three voltages, same \SI{3032}{\watt} array.
$A_{\min}$ is the conductor for a \SI{3}{\percent} drop over \SI{3}{\meter};
$P_{\text{loss}}$ is the round-trip $I^2R$ in a fixed \SI{16}{\milli\meter\squared}
cable.}
\label{tab:busvolt}
\small
\begin{tabular}{l S[table-format=3.0] S[table-format=3.0] S[table-format=2.2] S[table-format=3.1] S[table-format=2.2]}
\toprule
Bus & {$I_{\text{bus}}$ (\si{\ampere})} & {Tiles/string}
    & {$A_{\min}$ (\si{\milli\meter\squared})}
    & {$P_{\text{loss}}$ (\si{\watt})} & {Loss (\si{\percent})} \\
\midrule
\SI{12}{\volt} & 253 & 1 & 36.21 & 411.7 & 13.58 \\
\SI{24}{\volt} & 126 & 2 & 9.05  & 102.9 & 3.39 \\
\SI{48}{\volt} & 63  & 4 & 2.26  & 25.7  & 0.85 \\
\bottomrule
\end{tabular}
\end{table}

The $I_{\text{bus}}$ values in Table~\ref{tab:busvolt} include a small allowance
for converter inefficiency above the ideal $P/V$, which is why the \SI{12}{\volt}
row reads \SI{253}{\ampere} rather than the ideal \SI{228}{\ampere}.

\subsection{Recommendation: raise the dock TEC bus to 48\,V}

Moving the TEC bus from \SI{12}{\volt} to \SI{48}{\volt} cuts the bus current by
a factor of four, from \SI{253}{\ampere} to \SI{63}{\ampere}. The required copper
drops from \SI{36}{\milli\meter\squared} (welding-cable territory) to
\SI{2.3}{\milli\meter\squared} (ordinary hookup wire), and the cable loss falls
from \SI{14}{\percent} to under \SI{1}{\percent} of the bank power. The MPPT
controller is rated by current, so a \SI{63}{\ampere} bus needs roughly a quarter
of the controller current of a \SI{253}{\ampere} bus; in practice this collapses
three paralleled \SI{12}{\volt} controllers into one, which roughly halves the
MPPT cost that the buy-list identified as oversized. A \SI{24}{\volt} bus is a
reasonable middle option (current and loss both cut by about \num{4} relative to
\SI{12}{\volt}'s loss, copper down to \SI{9}{\milli\meter\squared}), but
\SI{48}{\volt} is recommended for the full \num{38}-tile dock because it keeps the
copper and the controller smallest while staying within standard safe
extra-low-voltage limits.

The recommendation applies only to the dock TEC bus. The \SI{12}{\volt} carriage
bus is kept unchanged: its peak current is \SI{7.7}{\ampere} and its runs are
under \SI{1.5}{\meter}, so there is no copper or loss penalty to fix, and
\SI{12}{\volt} keeps the motor, pump, and logic supplies simple.

%=====================================================================
\section{Step 5 --- Wire and fuse sizing}
\label{sec:wire}

Conductors are sized by Eq.~\eqref{eq:amin} and then checked against the
ampacity table with the IEC~60364-5-52 derating factor of \num{0.87} for
\SI{40}{\degreeCelsius} ambient. Fuses are sized at
\begin{equation}
  I_{\text{fuse}} = 1.25\,I_{\text{rated}},
  \label{eq:fuse}
\end{equation}
rounded up to the next standard rating, with the fuse hold current above the
load and the trip below the relevant stall or fault current.

\subsection{Carriage branches}

Table~\ref{tab:carriage-wire} sizes the carriage branches. The minimum
cross-sections are tiny because the currents and runs are small; the practical
choice is set by mechanical robustness and connector standards, not by the
electrical minimum. The carriage uses \SI{16}{AWG} (\SI{1.3}{\milli\meter\squared})
throughout, which clears every branch with large margin.

\begin{table}[H]
\centering
\caption{Carriage wire and fuse sizing at \SI{12}{\volt}, \SI{3}{\percent} drop,
\SI{0.87} derating at \SI{40}{\degreeCelsius}.}
\label{tab:carriage-wire}
\small
\begin{tabular}{l S[table-format=1.2] S[table-format=1.2] S[table-format=1.1] S[table-format=2.2] c}
\toprule
Branch & {$I$ (\si{\ampere})} & {Peak (\si{\ampere})} & {$L$ (\si{\meter})}
       & {$A_{\min}$ (\si{\milli\meter\squared})} & {Fuse} \\
\midrule
Pump            & 3.75 & 4.0 & 1.5 & 0.27 & \SI{5}{\ampere} \\
Drive (2 motors)& 1.61 & 3.0 & 1.0 & 0.16 & \SI{5}{\ampere} \\
Brush           & 0.13 & 1.6 & 1.0 & 0.15 & \SI{3}{\ampere} \\
Main feed       & 6.2  & 8.0 & 0.5 & 0.16 & \SI{10}{\ampere} \\
\bottomrule
\end{tabular}
\end{table}

The drive and brush peak columns are the start-up inrush near stall; the fuses
are chosen to hold the inrush but trip well below a hard short. The main-feed
\SI{10}{\ampere} fuse sits below the BMS \SI{15}{\ampere} limit so the fuse, not
the BMS, clears a carriage fault.

\subsection{Dock TEC bus at 48\,V}

At the recommended \SI{48}{\volt} bus the main run carries
$I_{\text{bus}}=\SI{63}{\ampere}$. From Eq.~\eqref{eq:amin} the
\SI{3}{\percent}-drop minimum is \SI{2.3}{\milli\meter\squared}, but ampacity
governs here, not voltage drop: after the \num{0.87} derate the conductor must
carry $\SI{63}{\ampere}/\num{0.87}=\SI{73}{\ampere}$, which calls for about
\SI{16}{\milli\meter\squared} (\SI{6}{AWG}) copper by the IEC ampacity table. The
main bus fuse is $1.25\times\SI{63}{\ampere}=\SI{79}{\ampere}\to\SI{80}{\ampere}$.

Each string carries the per-tile \SI{6}{\ampere}, so each of the ten strings
(nine quads and one pair) is individually fused at
$1.25\times\SI{6}{\ampere}=\SI{7.5}{\ampere}\to\SI{8}{\ampere}$ and switched by
its own MOSFET. The switch is an IRLB4030 logic-level N-MOSFET
\cite{irlb4030ds}: an open string can raise its node to the full string voltage,
so the part must stand off at least $2\times\SI{53}{\volt}=\SI{106}{\volt}$ and
carry at least $2\times\SI{6}{\ampere}=\SI{12}{\ampere}$, which the IRLB4030
(\SI{100}{\volt}, \SI{180}{\ampere}) clears on current with room and meets on
voltage. Per-string switching and fusing matter together: a single shorted tile
in a \num{38}-tile bank must blow its own \SI{8}{\ampere} fuse and drop only its
string without taking down the \SI{80}{\ampere} main, so the rest of the bank
keeps harvesting.

%=====================================================================
\section{Step 6 --- Protection}
\label{sec:protection}

Table~\ref{tab:protection} lists the protection elements for both subsystems,
following the workflow's selection criteria.

\begin{table}[H]
\centering
\caption{Protection elements, carriage and dock.}
\label{tab:protection}
\small
\begin{tabular}{p{3.0cm} p{4.3cm} p{6.0cm}}
\toprule
Hazard & Element & Selection \\
\midrule
Reverse polarity & P-channel MOSFET in the battery feed
  & $V_{DS}\geq 1.3\,V_{\text{supply}}$; low $R_{DS(on)}$ to minimise drop. A
    \SI{60}{\volt} P-MOSFET covers both the \SI{12}{\volt} carriage and the
    \SI{48}{\volt} dock bus. \\
Overcurrent & Blade fuses, per Table~\ref{tab:carriage-wire} and the dock bus
  & $I_{\text{fuse}}=1.25\,I_{\text{rated}}$; hold above load, trip below stall
    or fault. Per-string \SI{8}{\ampere} fuses isolate a single failed tile. \\
Motor transients & Flyback diode across each motor and switch; RC snubber across
  the H-bridge & Diode clamps the inductive kick at PWM switch-off; the snubber
    damps ringing on the MDD3A and the low-side MOSFET outputs. \\
Sensor / comm ESD & TVS diode on each I\textsuperscript{2}C, 1-Wire, and encoder
  line & $V_{\text{clamp}}\leq\num{0.9}\,V_{\text{breakdown}}$ of the protected
    input; \SI{5}{\volt} working voltage TVS on the \SI{3.3}{\volt}/\SI{5}{\volt}
    logic lines. \\
TEC bus overvoltage & Zener/TVS clamp across each string
  & Trigger $\leq V_{\text{abs,max}}$ of the tile (string $V$ stays below
    $4\times V_{\max}$). Protects against an open string raising the bus. \\
Pack abuse & LFP BMS (over-charge, over-discharge, over-current, balance)
  & \SI{15}{\ampere} continuous limit; the \SI{10}{\ampere} main fuse clears
    faults before the BMS. \\
\bottomrule
\end{tabular}
\end{table}

The reverse-polarity P-MOSFET is preferred over a series diode because it drops
only $I\,R_{DS(on)}$ (tens of millivolts) rather than a diode's \SI{0.7}{\volt},
which matters on the \SI{12}{\volt} carriage where every dropped volt is fuel
budget. Flyback diodes are mandatory on all three motors because the MDD3A and
the brush and pump MOSFETs switch inductive loads at \SI{20}{\kilo\hertz};
without them the switch-off spikes from Eq.~\eqref{eq:ripple}'s ripple current
would stress the driver and the low-side FETs.

%=====================================================================
\section{Committed parts and carriage mass closure}
\label{sec:weight}

\subsection{Committed parts}

Table~\ref{tab:committed} fixes one specific part per function with its governing
electrical numbers, so the buy-list and the drawings build from a single
selection rather than a menu of options.

\begin{table}[H]
\centering
\caption{Committed power and control parts. One part per function; the BTS7960
of the earlier draft is replaced by the right-sized MDD3A.}
\label{tab:committed}
\small
\begin{tabular}{p{3.1cm} p{4.4cm} p{5.6cm}}
\toprule
Function & Committed part & Key numbers \\
\midrule
Drive motor ($\times2$) & GA25-370 \SI{12}{\volt}, \num{1}:\num{71}, encoder
  & \SI{0.78}{\newton\meter} stall, \SI{60}{\rpm} no-load, \SI{1.5}{\ampere}
    stall, \SI{0.81}{\ampere} at load; \SI{0.105}{\kilo\gram} each. \\
Drive hold & MDD3A brake + sprung park detent (worm-gear alt.)
  & holds \SI{0.24}{\newton\meter\per rail} back-drive with supply off. \\
Motor driver & Cytron MDD3A, dual channel \cite{cytronmdd3a}
  & \SI{3}{\ampere} cont., \SI{5}{\ampere} peak, \SIrange{5}{30}{\volt}; one
    module, both rail motors; \SI{0.020}{\kilo\gram}. \\
Brush motor & GA25-370 \SI{12}{\volt}, \SI{100}{\rpm}
  & \SI{0.13}{\ampere} steady, \SI{1.6}{\ampere} inrush;
    \SI{0.105}{\kilo\gram}. \\
Brush + pump switch & IRLZ44N low-side MOSFET \cite{irlz44nds}
  & \SI{55}{\volt}, \SI{47}{\ampere}, logic-level gate; one per load. \\
Water pump & \SI{12}{\volt} diaphragm, \SIrange{3.6}{4.0}{\liter\per\minute}
  & \SI{3.75}{\ampere}, \SI{45}{\watt}; \SI{0.30}{\kilo\gram}. \\
Carriage battery & \SI{12}{\volt} \SI{10}{\ampere\hour} LFP, \SI{15}{\ampere} BMS
  & \SI{91}{\watt\hour} usable, $\sim$5 sessions; \SI{1.30}{\kilo\gram}. \\
TEC string switch ($\times10$) & IRLB4030 \SI{100}{\volt} MOSFET
  \cite{irlb4030ds} & one per series-string, \SI{6}{\ampere} string current. \\
Dock MPPT & EPever Tracer4215AN \cite{epever2023tracer42}
  & \SI{40}{\ampere}, \SI{150}{\volt} PV, \SI{48}{\volt} bus; one unit. \\
Real-time MCU & ESP32 DevKit \cite{esp32ds}
  & drive PWM, encoders, TEC PID; \SI{0.010}{\kilo\gram}. \\
ROS~2 host & Raspberry~Pi~4B \cite{rpi2021brief}
  & power-gated outside cleaning; \SI{0.055}{\kilo\gram}. \\
\bottomrule
\end{tabular}
\end{table}

\subsection{Carriage mass closure}

The CAD package froze the carriage at \SI{8.0}{\kilo\gram} with the TEC condenser
counted on the carriage and the battery buried inside a \SI{0.37}{\kilo\gram}
``electronics bay'' line. Two facts break that closure once the parts are
real. First, the design pivoted the TEC condenser to the stationary dock, so its
\SI{1.05}{\kilo\gram} no longer rides the carriage. Second, an honest
\SI{12}{\volt} \SI{20}{\ampere\hour} LFP pack weighs about
\SI{2.3}{\kilo\gram}, not \SI{0.37}{\kilo\gram}; keeping both the TEC and a
\SI{20}{\ampere\hour} pack on the carriage would put it near \SI{11}{\kilo\gram}.

The fix is the two corrections together: the TEC leaves the carriage with the
pivot, and the pack drops to \SI{12}{\volt} \SI{10}{\ampere\hour}
(\SI{1.3}{\kilo\gram}) with dock top-up charging in place of one large pack.
Table~\ref{tab:weight} shows the carriage then closes at
\SI{7.99}{\kilo\gram}, inside the \SI{8.0}{\kilo\gram} target.

\begin{table}[H]
\centering
\caption{Itemised carriage mass at mean tank fill. TEC condenser moved to the
dock; battery downsized to \SI{10}{\ampere\hour}. Closure \SI{7.99}{\kilo\gram}.}
\label{tab:weight}
\small
\begin{tabular}{l S[table-format=1.3]}
\toprule
Item & {Mass (\si{\kilo\gram})} \\
\midrule
Carriage base plate (Al 6061, \SI{3}{\milli\meter}) & 1.100 \\
Full-width brush \SI{0.96}{\meter}                   & 1.000 \\
Carriage side-frames ($2\times$)                     & 0.800 \\
V-groove guide wheels ($8\times$) + bogie plates     & 0.400 \\
Drive motors ($2\times$ GA25-370)                    & 0.210 \\
Steel pinions ($2\times$ $\varnothing$\SI{30}{\milli\meter} M1) & 0.060 \\
Brush motor (GA25-370)                               & 0.105 \\
Driver (MDD3A) + brush/pump MOSFETs                  & 0.050 \\
Water pump (diaphragm)                               & 0.300 \\
Battery (\SI{12}{\volt} \SI{10}{\ampere\hour} LFP)   & 1.300 \\
ESP32 DevKit                                         & 0.010 \\
Raspberry~Pi~4B + heatsink                           & 0.055 \\
Sensors + 2 limit switches                           & 0.100 \\
Cabling, connectors, fuses                           & 0.300 \\
Tank dry + spray bar + nozzles + tubing              & 0.400 \\
Water, mean fill (\SI{1.8}{\liter})                  & 1.800 \\
\midrule
\textbf{Total}                                       & \textbf{7.990} \\
\bottomrule
\end{tabular}
\end{table}

%=====================================================================
\section{Design summary and quality checks}
\label{sec:summary}

\subsection{Key results}

\begin{itemize}
\item \textbf{Carriage battery:} \SI{12}{\volt} \SI{10}{\ampere\hour} LFP,
  \SI{91}{\watt\hour} usable, giving about \textbf{5 full-array cleaning
  sessions per charge} (roughly \num{18} days at the \num{3.5}-day interval),
  with a dock top-up between runs. The pack is held to \SI{10}{\ampere\hour}
  rather than \SI{20}{\ampere\hour} to keep the carriage under its
  \SI{8}{\kilo\gram} limit (Section~\ref{sec:weight}). Peak draw
  \SI{6.2}{\ampere} (\SI{0.62}{C}); continuous runtime about \SI{1.8}{\hour}.
\item \textbf{Drive:} two GA25-370 \num{1}:\num{71} motors, \SI{0.81}{\ampere}
  each, one Cytron MDD3A dual driver (right-sized; the old BTS7960 was
  \num{50}$\times$ over-rated), \num{1.86}$\times$ torque margin per motor. The
  carriage holds on the tilt by the driver brake plus a park detent.
\item \textbf{Dock TEC bus:} \textbf{raise from \SI{12}{\volt} to
  \SI{48}{\volt}} (nine four-tile strings plus one pair). This cuts bus current
  from \SI{253}{\ampere} to \SI{63}{\ampere}, copper from
  \SI{36}{\milli\meter\squared} to \SI{2.3}{\milli\meter\squared} (ampacity sets
  \SI{16}{\milli\meter\squared}), and cable loss from \SI{14}{\percent} to under
  \SI{1}{\percent}, and collapses three \SI{12}{\volt} MPPT units to one
  \SI{48}{\volt} controller. The \SI{12}{\volt} carriage bus is kept as-is.
\item \textbf{TEC operating point:} \SI{6}{\ampere}, \SI{13.3}{\volt},
  \SI{79.8}{\watt} per tile, $Q_C=\SI{37.4}{\watt}$, instantaneous
  $COP=\num{0.47}$ (system night-strategy COP \num{1.9} is a separate seasonal
  figure).
\end{itemize}

\subsection{Quality checklist}

\begin{itemize}
\item[\checkmark] Power budget completed before any component was selected
  (Section~\ref{sec:budget}).
\item[\checkmark] Motor stall current verified against driver peak rating
  (Eq.~\eqref{eq:mdd3a}).
\item[\checkmark] Wire sizes derated \num{0.87} for \SI{40}{\degreeCelsius}
  ambient (Section~\ref{sec:wire}).
\item[\checkmark] Reverse-polarity, fuse, TVS, flyback, and BMS protection all
  specified (Table~\ref{tab:protection}).
\item[\checkmark] Battery capacity includes \SI{20}{\percent} margin and the LFP
  $DOD=\num{0.8}$ (Section~\ref{sec:storage}).
\item[\checkmark] ICs cited with manufacturer, part number, and year where a
  datasheet exists (TEC1-12706, Cytron MDD3A, IRLZ44N, IRLB4030, INA219, ESP32,
  Raspberry~Pi).
\item[\checkmark] Equations labelled; variables defined immediately after each.
\item[\checkmark] Operating temperature confirmed: LFP, TEC, and logic all rated
  for the \SI{40}{\degreeCelsius} NAC ambient.
\end{itemize}

%=====================================================================
\bibliographystyle{IEEEtran}
\bibliography{references}

\end{document}
